\cotefile{162-5}
\scanpage{41}
\unsure{\deleted{L}C}\unsure{\deleted{a}e}s\deleted{e} bicomplexes s'identifient aux complexes de c\unsure{\deleted{a}h}\unsure{\deleted{h}a}ines de lo\unsure{\deleted{g}n}\unsure{\deleted{n}g}ueur 2.\linebreak
Donc la catégorie de (2-Cat)-groupes abéliens est équivelente à celle des\linebreak
complexes de chaines de longueur 2. L'explicitation de la 2-catégorie\linebreak
associée à un tel complexe de chaines est mon interprétation chérie. Il\linebreak
faudrait encore expliciter l\unsure{\deleted{e}a}\deleted{s morphismes} 2-Grab-catégorie des\linebreak
homomorphismes (en un sens à préciser) entre deu\unsure{\deleted{t}x}\deleted{el} tel\unsure{\deleted{s}l}es 2-Grab-catégo-\linebreak
ries strictes, en termes d'un complexe $\mathrm{RHom}$.

L'idée qui se dégage de cette approche, et que Quillen on espère va\linebreak
expl\unsure{\deleted{o}i}citer, c'est que la \deleted{notion de catégorie de} notion\deleted{s} de groupe abélien\linebreak
dans (n-Cat) correspond à celle de complexe de \inserted{chaines de} longueur n, en raisonnant\linebreak
soit mod. isomorphisme de part et d'autre, soit mod. une notion naturelle\linebreak
de n-équivalence d'un côté, et dans la catégorie dérivée de l'autre. Ceci\linebreak
doit pouvoir se préciser en utilisant des $\mathrm{Hom}^{\bullet}$ resp. des $\mathrm{RHom}^{\bullet}$.

Si maintenant C est une catégorie avec un ``produit tensoriel'' associa\linebreak
tif (et éventuellement unitaire), i.e. un bifoncteur avec des isomorphis-\linebreak
mes d'associativité (mais sans inverse ...), on peut lui associer \unsure{\deleted{(}s}ans\linebreak
doute une Gr-catégorie universelle (tout comme à un monoïde on associe un\linebreak
groupe symmétrisé), d'où des invariants, qu'on pourra désigner par $K_o(C)$\linebreak
et $K_1(C)$. On suppose qu'on pourra de même trouver, pour tout n, une\linebreak
Gr-n-catégorie universelle où envoyer C de façon compatib\unsure{\deleted{o}l|\deleted{b}l}e avec $\otimes$ (\unsure{\deleted{d}t}out\linebreak
cela dans un sens à préciser), dont les invariants $\pi_n$ seraient par défini\linebreak
tion les $K_n(C)$. \unsure{\deleted{E}U}\unsure{\deleted{e}n} cas intéressant est celui où C est une catégorie avec\linebreak
produits finis (par exemple une catégorie additive) et $\otimes$ est le produit\unsure{;|,}\linebreak
\deleted{mais} \deleted{(}\inserted{[}Mais on notera que dans le cas d'une catégorie non seulement additive\linebreak
mais abélienne, il y a une notion plus adéquate du $K_o$, qui devrait corres\linebreak
pondre à une notion plus adéquate de la suite des $K_i$, inspirée du tapis\linebreak
de Roos; sans doute c'est en termes de (``vraie'') catégories triangulées

\margin{40}
